\chapter*{General introduction}
\addcontentsline{toc}{chapter}{General introduction}
\markboth{General introduction}{General introduction}
\label{chap.intro}

\section*{Context}

Warehouses sit at the centre of modern supply chains. The growth of e-commerce
and the demand for ever faster delivery have multiplied the number of references
stored, the size of the buildings that hold them, and the pace at which goods
move in and out. In this setting, knowing exactly what is on the shelves is not
an administrative detail. An inaccurate stock record leads to failed orders,
unnecessary purchases and capital tied up in goods that nobody can find.

Yet inventory remains one of the least automated operations of a warehouse. It is
still largely carried out by operators who walk the aisles with handheld scanners
and use lifts to reach the upper levels of the racks. The task is slow,
repetitive and exposed to human error, it takes staff away from other work, and
it must be repeated regularly for the records to stay reliable. As buildings grow
larger and taller, the cost of every inventory campaign grows with them.

Autonomous aerial vehicles offer a natural answer to this problem. A drone moves
freely in three dimensions, reaches the highest shelves without any equipment,
and carries a camera able to read the labels already present on the goods.
Commercial solutions and research prototypes have already shown drones scanning
racks in operating warehouses, and the idea of using several vehicles together,
to cover a whole building in a fraction of the time, has followed naturally.

This work is positioned within this context. It investigates how a team of
autonomous drones can carry out the inventory of a warehouse on its own, from
take-off to the delivery of a complete list of the goods the building holds, each
one located where it stands.

\section*{Problem Statement}

Despite this potential, autonomous inventory by a team of drones remains an open
challenge. Most existing approaches rely on a map of the warehouse provided in
advance: the route of each vehicle is computed before take-off, and the work is
divided between the vehicles once and for all. These systems perform well when
the mission unfolds as planned, but they become fragile as soon as it does not. A
vehicle that fails leaves its share of the work undone, an aisle blocked during
the mission invalidates the planned route, and a building whose map is missing or
outdated cannot be inventoried at all.

Removing this dependency on a prior map raises several difficulties at once. The
team must decide where to look without knowing where the goods are, and build its
knowledge of the building while exploring it. It must share the work among its
members without a central planner, which would become a single point of failure,
while avoiding duplicated effort and collisions in narrow aisles. And it must
keep delivering a complete inventory when the conditions of the mission change.

This work addresses the challenge of building a team of drones that inventories a
warehouse without any prior knowledge of its layout, and studies what makes such
a team complete, efficient and robust. Four research questions guide this work:

\begin{itemize}[itemsep=0.4em]
    \item \textbf{RQ1}: How can a team of drones decide where to go and what to
          read in a warehouse it does not know, using only the knowledge it
          gathers during the flight?
    \item \textbf{RQ2}: How can the vehicles share the work without a central
          planner, while avoiding duplicated effort and remaining safe in narrow
          aisles?
    \item \textbf{RQ3}: How can the team remain robust when the mission does not
          unfold as expected, such as the loss of a vehicle, an obstacle appearing
          in an aisle, or a warehouse never seen before?
    \item \textbf{RQ4}: How can a team that starts without any map outperform
          static approaches that plan the whole mission in advance from a known
          map?
\end{itemize}

\section*{Contributions}

To answer these questions, this work makes the following contributions:

\begin{itemize}[itemsep=0.4em]
    \item A decentralised inventory system for a team of drones that requires no
          map of the warehouse and builds a shared representation of the building
          during the flight.
    \item A perception chain that combines a learned detector with a
          \glsxtrshort{abb:qr} decoder, able to spot labels at a distance and to
          locate every identifier it reads.
    \item A semantic guidance module based on a \gls{abb:vlm}, which advises the
          vehicles on the areas to visit next.
    \item A simulation environment for multi-drone warehouse inventory, with
          procedurally generated warehouses and realistic flight control.
    \item An experimental evaluation of the robustness of the system, in nominal
          conditions and when a vehicle is lost, an obstacle appears in an aisle,
          or the warehouse has never been seen before.
\end{itemize}

\section*{Report Organization}

This report presents the work carried out as part of this final year project and
is organised into two parts.

The first part presents the background and the state of the art:

\begin{itemize}[itemsep=0.4em]
    \item \textbf{\Cref{chap:background}} introduces warehouse inventory and the
          use of aerial robots in logistics, and presents the limits of current
          practices.
    \item \textbf{\Cref{chap:nav}} reviews existing approaches for autonomous
          inventory, multi-robot exploration and task coordination, and identifies
          the gap addressed by this work.
\end{itemize}

The second part presents the proposed contribution:

\begin{itemize}[itemsep=0.4em]
    \item \textbf{\Cref{chap:system}} presents the design of the proposed system,
          its overall architecture and its main components.
    \item \textbf{\Cref{chap:implementation}} describes the implementation of the
          system and its experimental evaluation.
\end{itemize}

Finally, the report concludes with a general conclusion that summarises the
contributions and the results obtained, and outlines the directions for extending
this work.
